% !TeX root = ../papers/nc-p3.tex
\section{二次元単純の固有な極限と有限窓の分岐}
\label{sec:finite-two-dimensional-specialization}
\begin{theorem}\label{thm:finite-two-dimensional-specialization}
十五次元混合成分の閉包の標準有限Frobenius窓で四辺が有限とする。
各頂点二次元の単純が存在しないなら、長さ二十の点グラフに次のいずれかがある。
\begin{enumerate}
\item 非被約な局所点。
\item 異なる二点 $p,q$ で $\operatorname{rank} w(p_{i+3},-,q_{i+1},-)<2$。
\end{enumerate}
とくに全点が被約でこの異点階数が常に2なら、その閉包の\emph{全ての}窓が
定理\ref{thm:finite-two-dimensional-quadrics}の二次曲面表示を少なくとも一つ持つ。
閉点のAS性や中心の有限性は仮定しない。
より一般に、この結論は一般点に二次元単純がある任意の既約な有限窓の族に通る。
\end{theorem}
\begin{proof}
連続四頂点の大切断 $B$ の表現を、矢が $0,1,2,3$ の順に作用する向きで読む。
$\delta=(1,1,1,1)$ とし、重み $\theta=(1,3,9,-13)$ について非零真部分表現の重みが
負となることを安定、零以下を半安定とする。$0\le d\le2\delta$ で
$\theta d=0$ の可能性は $0,\delta,2\delta$ だけである。三進表示で最初の三成分が決まるためである。
四辺で前節の全射条件を満たす $2\delta$ 表現の切断は安定である。
非零部分が途中から始まっても一回作用させれば次位置は全体となる。
最後の位置の次元が2なら前三位置の真の欠損で重みは負となり、
最後の次元1の部分は最後にのみ支えられるからやはり負である。
安定な $\delta$ 表現は薄い点 $M(p)$ に限る。途中の矢が零なら正重みの初期区間が部分表現となり、
全て非零なら有限な点列である。

逆に切断で安定な $2\delta$ 表現の三辺では、各非零ベクトル上の矢の像が二次元となる。
正重みの各頂点には共通核がない。第一辺で一つの像が直線なら、有限辺の評価階数3以上により
その後も直線に入り、初期から末尾までの $\delta$ 部分または正重みの初期部分を生じる。
双対を逆に読んでも初期区間の重み和は真区間で正（元の終期区間の符号反転）なので、
最終辺でも双対の像は全非零ベクトル上で二次元である。
二行二列の場合この条件は元の像条件と同値である。すなわち共変元とベクトルの非零対が
全矢を同時に消すことを排除するだけである。第二辺で直線が生じれば、
有限辺の関係により最終辺でも直線以下になり矛盾する。
これで安定な切断の三辺のいずれも像条件を満たす。

第四辺を付ける全ての係数に対し、残る二窓の関係は斉次線型である。
安定な切断で非零解が一つでもあれば、それが与える折り畳み表現は単純である。
実際、第四辺であるベクトル上の像が直線なら第一辺へ伝播して矛盾する。
像が零のベクトルがあれば、第四辺の全矢は共通の一次元商を通る。
第四辺は非零なので双対側ではある共変元上の像が直線になり、逆向きに最終辺へ伝播して矛盾する。
従って四辺で全射条件が成立する。

一般窓からの離散付値環上の特殊化を取る。非巡回な四頂点箙に関係を課した
$2\delta$ の $\theta$ 半安定表現の相対GIT商は底に射影的である。
これは半不変式環のProjであり、次数零の不変式は有向閉路がないため底の関数だけである。
半安定縮約により有限な体拡大後、一般の安定表現は特殊繊維が半安定な表現族を持つ
（箙のGITの半安定置換、またはQuot表示による半安定格子の選択）。
関係式の閉条件は保存される。第四辺の解空間が非零という条件も行列階数の閉条件なので極限に残る。
従って特殊な切断が安定なら既に結論を得る。

安定でなければ安定因子は $M(p),M(q)$ の二つである。
$p=q$ を含め、両局所点が単純で異点の反対縮約階数が2なら、
定理\ref{thm:finite-cut-point-pairings}から二因子間・自己の第一Extは全て零である。
従って特殊表現は直和 $N=M(p)\oplus M(q)$ で $\operatorname{Ext}^1_B(N,N)=0$。
長さ二十の有限完全交叉はこの族で平坦で、単純な点はエタール局所で切断へ延びる。
直和の変形と選んだ半安定格子の変形は、完備化後、各次数で基底変更を持ち上げて同型になる。
実際隣り合う冪零厚みで二つの持上げの差は自己第一Extで測られ、零なら
前段の同型を修正して延長できる。行列の完備極限で同型を得るので一般点も直和となり安定性に反する。
一般点の存在は前節までの中心Euler数の結果による。
任意の閉包点への移行には稠密な存在開集合を一般点とする曲線の付値環を使えばよい。
\end{proof}

\begin{corollary}\label{cor:two-dimensional-extension-space}
上の安定な切断 $N$ に対する第四辺の線型解空間は自然に
$D\operatorname{Ext}^3_B(N,N)$ と同型である。従って追加の単純の存在と、
この安定表現が最高自己Extを持つことは同値である。
\end{corollary}
\begin{proof}
戻り辺を付ける表示は $T_B(\operatorname{Ext}^3_B(DB,B))$ なので、
解空間は $\operatorname{Hom}_B(N\otimes_B \operatorname{Ext}^3_B(DB,B),N)$ である。
大域次元3以下、命題\ref{prop:two-window-ext-complex}とSerre関手 $S$ の式より、
左の引数は $H^0(S^{-1}N[3])$ で、負次数以外に他のコホモロジーはない。
次数零加群への導来射はこの $H^0$ からの射と一致する。Serre双対で
$\operatorname{Hom}(S^{-1}N[3],N)=D\operatorname{Hom}(N,N[3])$ となる。
\end{proof}
